\documentclass[10pt, letterpaper]{article}

% Inhaltsverzeichnis für Pakettypen (nur für Übersicht im Header, wird nicht im Dokument angezeigt)
% 1. Seitenlayout und Ränder
% 2. Sprache und Zeichensatz
% 3. Mathematik und Theorem-Umgebungen
% 4. Eigene Makros
% 5. Diagramme und Grafiken
% 6. Tabellen und Aufzählungen
% 7. Inhaltsverzeichnis
% 8. Abschnittsüberschriften
% 9. Abstrakt-Umgebung
% 10. Todos/Notizen
% 11. Rahmen/Box-Umgebungen
% 12. Python-Integration
% 13. Literaturverwaltung
% 14. Hyperlinks
% 15. Absatzeinstellungen
% 16. Umgebungen
% 17  Graphik
% 18  Extra
% 00. Titel und Autor

\usepackage{slashed}

% --- 1. Seitenlayout und Ränder ---
\usepackage[margin=3cm]{geometry}

% --- 2. Sprache und Zeichensatz ---
\usepackage[english]{babel}
\usepackage[T1]{fontenc}
\usepackage[utf8]{inputenc}

% --- 3. Mathematik und Theorem-Umgebungen ---
\usepackage{amsmath, amssymb, amsthm}
\usepackage{mathrsfs}
\DeclareMathOperator{\WF}{WF}

% --- 4. Eigene Makros ---
\usepackage{xcolor}
\newcommand{\SKP}{\langle\cdot,\cdot\rangle}
\newcommand{\R}{\mathbb{R}}
\newcommand{\N}{\mathbb{N}}
\newcommand{\Q}{\mathbb{Q}}
\newcommand{\Z}{\mathbb{Z}}
\newcommand{\C}{\mathbb{C}}
\newcommand{\entwurf}[1]{\textcolor{red}{#1}}
\newcommand{\GL}{\mathrm{GL}}
\newcommand{\SL}{\mathrm{SL}}
\newcommand{\SO}{\mathrm{SO}}
\newcommand{\SU}{\mathrm{SU}}
\newcommand{\Sp}{\mathrm{Sp}}
\newcommand{\Ogroup}{\mathrm{O}}
\newcommand{\U}{\mathrm{U}}

% --- 5. Diagramme und Grafiken ---
\usepackage{graphicx}
\graphicspath{{../../Library/Images/}}
\usepackage[export]{adjustbox}
\usepackage{tikz}
\usetikzlibrary{decorations.pathreplacing, arrows.meta, positioning}
\usepackage{tikz-cd}

% --- 6. Tabellen und Aufzählungen ---
\usepackage{enumitem}
\setlist[itemize]{left=0.5cm}

\newenvironment{romanenum}[1][]
  {%
    \ifx&#1&
    \else
      \textbf{#1}\quad
    \fi
    \begin{enumerate}[label=\roman*)]
  }
  {%
    \end{enumerate}%
  }

% --- 7. Inhaltsverzeichnis ---
\usepackage{tocloft}
\renewcommand{\cftsecfont}{\footnotesize}
\renewcommand{\cftsubsecfont}{\footnotesize}
\renewcommand{\cftsubsubsecfont}{\footnotesize}
\renewcommand{\cftsecpagefont}{\footnotesize}
\renewcommand{\cftsubsecpagefont}{\footnotesize}
\renewcommand{\cftsubsubsecpagefont}{\footnotesize}
\usepackage{etoc}

% --- 8. Abschnittsüberschriften ---
\usepackage{titlesec}
\titleformat{\section}{\normalfont\large\bfseries}{\thesection}{1em}{}
\titleformat{\subsection}{\normalfont\normalsize\bfseries}{\thesubsection}{0.5em}{}
\titleformat{\subsubsection}{\normalfont\normalsize\bfseries}{\thesubsubsection}{0.5em}{}
\setcounter{secnumdepth}{4}

% --- 9. Abstrakt-Umgebung ---
\usepackage{changepage}
\renewenvironment{abstract}
  {
    \begin{adjustwidth}{1.5cm}{1.5cm}
    \small
    \textsc{Abstract. –}%
  }
  {
    \end{adjustwidth}
  }

% --- 10. Todos/Notizen ---
\usepackage{todonotes}
% Hellblau definieren
\definecolor{hellblau}{rgb}{0.3,0.6,1}
\newenvironment{note}{\par\begingroup\color{hellblau}\itshape}{\par\endgroup}

% --- 11. Rahmen/Box-Umgebungen ---
\usepackage{mdframed}
\usepackage{tcolorbox}
\colorlet{shadecolor}{gray!25}

\newenvironment{customTheorem}
  {\vspace{10pt}%
   \begin{mdframed}[
     backgroundcolor=gray!20,
     linewidth=0pt,
     innertopmargin=10pt,
     innerbottommargin=10pt,
     skipabove=\dimexpr\topsep+\ht\strutbox\relax,
     skipbelow=\topsep,
   ]}
  {\end{mdframed}
   \vspace{10pt}%
  }

% --- 12. Python-Integration ---
% (Deaktiviert in dieser Version, aktiviere bei Bedarf)
% \usepackage{pythontex}
% \usepackage[makestderr]{pythontex}

% --- 13. Literaturverwaltung ---
\usepackage{csquotes}
\usepackage[backend=biber, style=alphabetic, citestyle=alphabetic]{biblatex}
\addbibresource{bibliography.bib}

% --- 14. Hyperlinks ---
\usepackage{hyperref}
\hypersetup{
  colorlinks   = true,
  urlcolor     = blue,
  linkcolor    = blue,
  citecolor    = blue,
  frenchlinks  = true
}

% --- 15. Absatzeinstellungen ---
\usepackage[parfill]{parskip}
\sloppy

% --- 16. Umgebungen ---
\usepackage{thmtools}

\newcommand{\CustomHeading}[3]{%
  \par\medskip\noindent%
  \textbf{#1 #2} \textnormal{(#3)}.\enskip%
}

\newenvironment{DEF}[2]{\begin{unitbox}\CustomHeading{Definition}{#1}{#2}}{\end{unitbox}}
\newenvironment{PROP}[2]{\begin{unitbox}\CustomHeading{Proposition}{#1}{#2}}{\end{unitbox}}
\newenvironment{THEO}[2]{\begin{unitbox}\CustomHeading{Theorem}{#1}{#2}}{\end{unitbox}}
\newenvironment{LEM}[2]{\begin{unitbox}\CustomHeading{Lemma}{#1}{#2}}{\end{unitbox}}
\newenvironment{KORO}[2]{\begin{unitbox}\CustomHeading{Corollar}{#1}{#2}}{\end{unitbox}}
\newenvironment{REM}[2]{\begin{unitbox}\CustomHeading{Remark}{#1}{#2}}{\end{unitbox}}
\newenvironment{EXA}[2]{\begin{unitbox}\CustomHeading{Example}{#1}{#2}}{\end{unitbox}}
\newenvironment{STUD}[2]{\begin{unitbox}\CustomHeading{Study}{#1}{#2}}{\end{unitbox}}
\newenvironment{CONC}[2]{\begin{unitbox}\CustomHeading{Concept}{#1}{#2}}{\end{unitbox}}
\newenvironment{OTH}[2]{\begin{unitbox}\CustomHeading{Other}{#1}{#2}}{\end{unitbox}}
\newenvironment{EXE}[2]{\begin{unitbox}\CustomHeading{Exercise}{#1}{#2}}{\end{unitbox}}
\newenvironment{MOT}[2]{\begin{unitbox}\CustomHeading{Motivation}{#1}{#2}}{\end{unitbox}}
\newenvironment{PROOF}[2]{\begin{unitbox}\CustomHeading{Proof}{#1}{#2}}{\end{unitbox}}

% --- Unit Umgebung für Source-Inhalte ---
\usepackage{mdframed}
\newmdenv[
  linewidth=1pt,
  topline=false,
  bottomline=false,
  rightline=false,
  leftmargin=0cm,
  rightmargin=0cm,
  skipabove=10pt,
  skipbelow=10pt,
  innertopmargin=0.5\baselineskip,
  innerbottommargin=0.5\baselineskip,
  backgroundcolor=gray!10,
  linecolor=gray
]{unitbox}

\newenvironment{unit}[1]
  {\begin{unitbox}\textbf{Unit #1}\par\smallskip}
  {\end{unitbox}}

% --- 17. ... ---

% --- 18. Extras ---
\usepackage{stmaryrd}
\usepackage{bbold}  % falls du \mathbb{1} nutzen willst

% --- 00. Titel und Autor ---
\title{Mein Titel}
\author{Tim Jaschik}
\date{\today}

\begin{document}

\maketitle
\rule{\textwidth}{0.5pt}
\begin{abstract}
Kurze Beschreibung …
\end{abstract}
\rule{\textwidth}{0.5pt}
\vspace{0.5cm}

\tableofcontents

\pagebreak


\section{Zeitabhängige Gauge-Gruppen}


\subsection{(Ko-)Homologische Struktur zeitabhängiger Eichgruppen}

\paragraph{(a) Ausgangspunkt: dynamische Eichstruktur.}
Sei $M$ eine glatte Raumzeitmannigfaltigkeit und
$G_\mathrm{gauge}(t)\subseteq G$ eine glatte Familie von Lie-Untergruppen.
Die zeitabhängige Eichgruppe definiert eine Familie von
Prinzipalbündeln $(P_t,\pi_t,M,G_\mathrm{gauge}(t))$
mit lokalen Verbindungen
\[
A_t(t)\in\Omega^1(M,\mathfrak g_\mathrm{gauge}(t)).
\]
Diese Verbindungen beschreiben lokale Zuordnungen
\[
A_t(t)\colon T_xM\to\mathfrak g_\mathrm{gauge}(t)
\]
und koppeln die innere Symmetrie an die Raumzeitstruktur.
Die Gesamtheit der glatten Verbindungen
\[
\mathcal A:=\{\,A\in\Omega^1(M,\mathfrak g)\,\}
\]
bildet eine affine Mannigfaltigkeit, auf der die Eichgruppe
\[
\mathcal G:=C^\infty(M,G)
\]
durch
\[
A\mapsto A^g:=g^{-1}Ag+g^{-1}dg
\]
wirkt.

\paragraph{(b) Kohomologische Beschreibung.}
Die Eichtransformationen $g\in\mathcal G$ sind
\emph{nichtlineare $0$-Kozyle} in einer nichtabelschen Kohomologie.
Der Zusammenhang $A$ definiert einen \emph{1-Kozyklus}
\[
A\in Z^1(M,\mathfrak g),
\qquad
dA+\tfrac12[A,A]=F\in Z^2(M,\mathfrak g),
\]
wobei $F$ die Krümmung des Zusammenhangs ist.
Zwei Verbindungen, die durch eine Eichtransformation verknüpft sind,
gehören zur selben Kohomologieklasse
\[
[A]\in H^1(M,G),
\]
der Menge der Isomorphieklassen flacher $G$-Prinzipalbündel.

\paragraph{(c) Zeitabhängige Familie von Kohomologieklassen.}
Für eine dynamische Eichgruppe $G_\mathrm{gauge}(t)$
entsteht eine glatte Familie von Klassen
\[
[A_t(t)]\in H^1(M,G_\mathrm{gauge}(t)).
\]
Ihre zeitliche Variation wird durch die zusätzliche
Verbindungskomponente $A_t(t)\,dt$
auf dem Produkt $M\times\mathbb R_t$ beschrieben.
Die Gesamtverbindung lautet
\[
\mathbb A = A(x,t) + A_t(x,t)\,dt
\;\in\;
\Omega^1(M\times\mathbb R_t,\mathfrak g),
\]
und die zugehörige Krümmung
\[
\mathbb F = d_{M\times\mathbb R}\mathbb A + \tfrac12[\mathbb A,\mathbb A]
= F_M + (D_MA_t-\partial_tA)\wedge dt.
\]
Das $dt$-Glied misst die zeitliche Veränderung der Kohomologieklasse:
\[
D_MA_t-\partial_tA=0
\quad\Leftrightarrow\quad
[A_t(t)]\text{ konstant in }H^1(M,G).
\]
Bei Nichtverschwindung definiert es ein nichttriviales
Element in $H^2(M\times\mathbb R_t,\mathfrak g)$
— die topologische Signatur einer zeitabhängigen
Verformung der Eichstruktur.

\paragraph{(d) Homologische Perspektive.}
Homologisch interpretiert beschreibt der Pfad
$t\mapsto [A_t(t)]$ einen Zyklus in der Kohomologiemenge $H^1(M,G)$.
Ein nichttrivialer Rand
\[
\frac{d}{dt}[A_t]\neq0
\]
entspricht einer \emph{topologischen Transition}:
das System wechselt zwischen nichtisomorphen
Bündelklassen. In der klassischen Feldtheorie entspricht dies
z.\,B.\ der Entstehung oder Auslöschung von Monopolstrukturen,
in der Quantenmechanik der Akkumulation geometrischer Phasen.

\paragraph{(e) Physikalische Interpretation.}
\begin{itemize}
\item \textbf{Statische Kohomologie:}
  $[A_t]$ bleibt konstant $\Rightarrow$ topologische Stabilität der Symmetrie.
\item \textbf{Zeitabhängige Kohomologie:}
  $[A_t]$ variiert $\Rightarrow$ topologische oder anomale Entwicklung.
\item \textbf{Berry-Krümmung:}
  Für einen Parameterraum $\mathcal P$ ist die $U(1)$-Verbindung
  $A=i\langle n|\mathrm d n\rangle$
  eine $1$-Form auf $\mathcal P$ mit Krümmung
  $F=\mathrm dA\in H^2(\mathcal P,U(1))$.
  Ihre Kohomologieklasse ist die erste Chern-Klasse des
  Eigenraum-Bündels.
\item \textbf{Chern–Simons- und Wess–Zumino-Terme:}
  In $3$- bzw. $4$-Dimensionen definiert $\mathrm{Tr}(A\,dA+\tfrac23A^3)$
  eine Klasse in $H^3(M,\mathbb R)$, die globale Anomalien charakterisiert.
\end{itemize}

\paragraph{(f) Übersicht der ko-/homologischen Ebenen.}
\[
\boxed{
\begin{array}{lll}
\textbf{Kohomologieebene} & \textbf{Repräsentanten} & \textbf{Physikalische Bedeutung}\\[3pt]
\hline
H^0(M,G) & g(x) & \text{Eichtransformationen (0-Kozyle)}\\[3pt]
H^1(M,G) & [A] & \text{Bündel- bzw. Eichklassen}\\[3pt]
H^2(M,\mathfrak g) & [F] & \text{Feldstärken, Krümmungen}\\[3pt]
H^2(M,U(1)) & [F_B] & \text{Berry-Phase, Chern-Klasse}\\[3pt]
H^3(M,\mathbb Z) & [\mathrm{CS}] & \text{Anomalien, topologische Invarianz}\\[3pt]
\end{array}
}
\]

\paragraph{(g) Geometrisches Diagramm.}
\[
\begin{tikzcd}[row sep=large, column sep=large]
  & \Omega^1(M,\mathfrak g)
      \arrow[d, shift left=.4ex, "{\,d+\tfrac12[\ ,\ ]\,}"]
      \arrow[d, shift right=.4ex, swap, "{F}"]
      \arrow[r, "{\pi}"]
    & H^1(M,G)
      \arrow[d, "{\partial_t}"] \\
  & \Omega^2(M,\mathfrak g)
      \arrow[r, "{\mathrm{cl}}"]
    & H^2(M,\mathfrak g)
\end{tikzcd}
\]
Zeitabhängige Eichgruppen führen zu einem Fluss
$t\mapsto[A_t(t)]$ in $H^1(M,G)$.
Ein nichttrivialer Rand $\partial_t[A_t]\neq0$
projiziert auf eine Klasse $[F_{Mt}]\in H^2(M\times\mathbb R_t,\mathfrak g)$,
die die zeitliche Krümmung der Symmetriestruktur beschreibt.

\paragraph{(h) Zusammenfassung.}
\[
\boxed{
\begin{aligned}
&A\in\Omega^1(M,\mathfrak g)\;\Rightarrow\;[A]\in H^1(M,G),\\[3pt]
&F=dA+\tfrac12[A,A]\;\Rightarrow\;[F]\in H^2(M,\mathfrak g),\\[3pt]
&A_t(t)\Rightarrow [A_t(t)]\in H^1(M,G_\mathrm{gauge}(t)),\\[3pt]
&\partial_t [A_t]=0 \;\Leftrightarrow\;\text{Topologische Stabilität},\\[3pt]
&\partial_t [A_t]\neq0 \;\Leftrightarrow\;\text{Anomalie, Phasenübergang oder geometrische Phase.}
\end{aligned}
}
\]

\paragraph{(i) Interpretationshinweis.}
Kohomologisch beschreibt die Zeitentwicklung in einem System mit dynamischer Eichgruppe
nicht nur die Bewegung innerhalb eines Orbits, sondern auch
die \emph{zeitliche Veränderung der Orbitstruktur selbst}.
Der Zeitfluss $t\mapsto [A_t(t)]$ ist eine Kurve in der nichtabelschen
Kohomologie $H^1(M,G)$;
sein Rand in $H^2(M,\mathfrak g)$ misst, ob die
Eichstruktur topologisch invariant (flach) oder deformiert (anomale Entwicklung)
ist.
Die (ko-)homologischen Invarianten sind somit die
\emph{integralen Erhaltungssätze der Symmetriegeometrie}.


\subsection{Integrale und kohomologische Erhaltungssätze}

\paragraph{(a) Von der lokalen Krümmung zur globalen Invarianz.}
Die Krümmung $F=dA+\tfrac12[A,A]\in\Omega^2(M,\mathfrak g)$
ist lokal eine Differentialform, global jedoch
nur bis auf Eichtransformation bestimmt.
Ihre Invarianz zeigt sich in der Kohomologieklasse $[F]\in H^2(M,\mathfrak g)$.

Ist $M$ orientiert und kompakt, so liefert das Integral über
geschlossene 2-Zyklen $C_2\subset M$
\[
\int_{C_2} \mathrm{Tr}(F)
\]
eine homologische Invariante:
das Integral hängt nur von der Kohomologieklasse $[F]$ und der Homologieklasse $[C_2]$
ab, nicht von der konkreten Wahl von $A$.
Dies ist die geometrische Form eines \emph{Erhaltungssatzes}.

\paragraph{(b) Chern-Klassen und quantisierte Flüsse.}
Für eine unitäre Darstellung $G\subset U(N)$
definiert $\mathrm{Tr}(F\wedge F)$ die \emph{zweite Chern-Klasse}
\[
c_2(A)=\frac{1}{8\pi^2}\mathrm{Tr}(F\wedge F)\in H^4(M,\mathbb Z),
\]
deren Integrale über geschlossene 4-Zyklen ganzzahlig sind:
\[
\int_{C_4} c_2(A)\in\mathbb Z.
\]
Diese Ganzzahligkeit entspricht der \emph{Quantisierung topologischer Ladungen}
(z.\,B.\ instanton number in Yang–Mills-Theorie).

\paragraph{(c) Chern–Simons-Invarianten und Randterme.}
Auf einer 3-dimensionalen Mannigfaltigkeit $M^3$
ist der Chern–Simons-Term
\[
\mathrm{CS}(A)=\mathrm{Tr}\!\left(A\wedge dA+\tfrac23A\wedge A\wedge A\right)
\]
eine $3$-Form mit
\[
d\,\mathrm{CS}(A)=\mathrm{Tr}(F\wedge F).
\]
Damit gilt:
\[
\int_{M^3}\mathrm{CS}(A)
\]
ist modulo $2\pi\mathbb Z$ wohldefiniert
und beschreibt eine \emph{topologische Phase} des Systems (z.\,B.\ Berry-Phase oder WZ-Term).

\paragraph{(d) Zeitentwicklung und integrale Erhaltung.}
Auf $M\times\mathbb R_t$ kann man
\[
\frac{d}{dt}\int_{C_2}\mathrm{Tr}(F)=\int_{\partial(C_2\times[0,t])}\mathrm{Tr}(F_{Mt})\,dt
\]
als Differentialform der Erhaltungssätze interpretieren.
Verschwindet $[F_{Mt}]$ in $H^2(M,\mathfrak g)$,
so ist der Fluss topologisch konservativ.

\paragraph{(e) Zusammenfassung.}
\[
\boxed{
\begin{aligned}
F\in Z^2(M,\mathfrak g)
&\Rightarrow [F]\in H^2(M,\mathfrak g),\\[3pt]
d\,\mathrm{CS}(A)&=\mathrm{Tr}(F\wedge F),\\[3pt]
\int_{C_4}\mathrm{Tr}(F\wedge F)\in 8\pi^2\mathbb Z,\\[3pt]
\partial_t[F]=0 &\;\Rightarrow\; \text{Erhaltung topologischer Ladung},\\[3pt]
\partial_t[F]\neq0 &\;\Rightarrow\; \text{Phasenübergang oder Anomalie.}
\end{aligned}
}
\]

\paragraph{(f) Interpretationshinweis.}
Diese integralen Größen sind die \emph{ko\-homologischen Schatten}
der dynamischen Symmetriestruktur:
\begin{itemize}
\item $H^2(M,\mathfrak g)$ misst die interne Krümmung (Feldstärke).
\item $H^3(M,\mathbb Z)$ klassifiziert topologische Phasen und Anomalien.
\item $H^4(M,\mathbb Z)$ kodiert quantisierte Instanton-Zahlen.
\end{itemize}
Physikalische Erhaltungssätze werden dadurch als
\emph{Integralklassen} verstanden, die sich aus der geschlossenen Struktur
der Krümmungsformen ergeben.








\pagebreak











\pagebreak
\printbibliography
\end{document}